%=====================================================================
%  DATABASE TECHNOLOGY, MANAGEMENT AND SECURITY
%  Project Type 8 --- Research Proposal / Mini Research Paper
%
%  Built on the course-provided template (db-report-template).
%  Compiler: pdfLaTeX. Bibliography: biber (IEEE style).
%=====================================================================

\documentclass[12pt, a4paper]{article}

%---------------------------------------------------------------------
% Packages
%---------------------------------------------------------------------
\usepackage[a4paper, margin=2.5cm]{geometry}
\usepackage[utf8]{inputenc}
\usepackage[T1]{fontenc}
\usepackage{times}
\usepackage{graphicx}
\usepackage{booktabs}
\usepackage{array}
\usepackage{multirow}
\usepackage{amsmath, amssymb}
\usepackage{caption}
\usepackage{enumitem}
\usepackage{setspace}
\usepackage{fancyhdr}
\usepackage{listings}
\usepackage{xcolor}
\usepackage{csquotes}
\usepackage[hidelinks]{hyperref}
\usepackage[nameinlink]{cleveref}

\usepackage[backend=biber, style=ieee, sorting=none]{biblatex}
\addbibresource{references.bib}

%---------------------------------------------------------------------
% Metadata
%---------------------------------------------------------------------
\newcommand{\CourseName}{Database Technology, Management and Security}
\newcommand{\ReportTitle}{LLM-Assisted Natural Language to SQL Query
Optimization Using Cost-Aware Database Planning}
\newcommand{\ProjectType}{Project Type 8 --- Research Proposal / Mini Research Paper}
\newcommand{\GroupNumber}{Individual Submission}
\newcommand{\SubmissionDate}{30 August 2026}
\newcommand{\InstructorName}{Dr.\ Ashraf Uddin}
\newcommand{\InstructorTitle}{Assistant Professor}
\newcommand{\InstructorDept}{Department of Computer Science}
\newcommand{\InstructorInstitution}{American International University-Bangladesh}

\newcommand{\ves}{\textsc{ves}}
\newcommand{\ex}{\textsc{ex}}

%---------------------------------------------------------------------
% Listing style
%---------------------------------------------------------------------
\definecolor{codegray}{gray}{0.95}
\definecolor{codegreen}{rgb}{0,0.5,0}
\definecolor{codeblue}{rgb}{0,0,0.7}
\lstdefinestyle{codestyle}{
    backgroundcolor=\color{codegray},
    basicstyle=\ttfamily\footnotesize,
    keywordstyle=\color{codeblue}\bfseries,
    commentstyle=\color{codegreen}\itshape,
    numbers=none, frame=single, framerule=0pt,
    breaklines=true, showstringspaces=false, captionpos=b, tabsize=2
}
\lstset{style=codestyle}

%---------------------------------------------------------------------
% Headers and footers
%---------------------------------------------------------------------
\pagestyle{fancy}
\fancyhf{}
\fancyhead[L]{\small\CourseName}
\fancyhead[R]{\small\GroupNumber}
\fancyfoot[C]{\thepage}
\renewcommand{\headrulewidth}{0.4pt}

\onehalfspacing
\emergencystretch=4em

\begin{document}

%---------------------------------------------------------------------
% TITLE PAGE
%---------------------------------------------------------------------
\begin{titlepage}
    \centering
    \vspace*{1.2cm}

    {\large \textbf{\CourseName}}\\[0.4cm]
    {\large \ProjectType}\\[1.6cm]

    \rule{\linewidth}{0.5mm}\\[0.6cm]
    {\LARGE \textbf{\ReportTitle}}\\[0.4cm]
    \rule{\linewidth}{0.5mm}\\[1.4cm]

    {\large \textbf{Submitted By}}\\[0.4cm]
    {\large \textbf{\GroupNumber}}\\[0.5cm]

    \begin{tabular}{@{}lll@{}}
        \toprule
        \textbf{Name} & \textbf{Student ID} & \textbf{Email} \\
        \midrule
        Md. Imtiaj Alam Sajin & 26-94090-2 & \small 26-94090-2@student.aiub.edu \\
        \bottomrule
    \end{tabular}\\[1.4cm]

    {\large \textbf{Submitted To}}\\[0.4cm]
    {\large \InstructorName}\\[0.15cm]
    \InstructorTitle\\[0.15cm]
    \InstructorDept\\[0.15cm]
    \InstructorInstitution\\[0.9cm]

    \vfill
    {\large Submission Date: \SubmissionDate}\\[0.4cm]
\end{titlepage}

%---------------------------------------------------------------------
% DECLARATION
%---------------------------------------------------------------------
\section*{Declaration of Originality}
\addcontentsline{toc}{section}{Declaration of Originality}
I declare that this proposal is my own original work. All sources of
information have been acknowledged and cited using the IEEE referencing
style.

\vspace{0.6em}
\noindent
\textbf{A note on the status of the results reported here.} This document
is a research \emph{proposal}, submitted under Project Type 8. It
therefore contains no experimental results for the system it proposes,
because that system has not yet been built. Every quantitative figure
that does appear, all of it confined to Section~\ref{sec:prelim}, comes
from a completed measurement study I carried out on PostgreSQL~17.1 and
MariaDB~10.4.28 during this course, and is reported there only because
it establishes the premise on which the proposal rests. No number in
this document is estimated, projected or illustrative. Where an outcome
is anticipated rather than measured, it is stated as an expectation and
labelled as such.

\vspace{1.2cm}
\noindent
\begin{tabular}{@{}p{6cm}@{}}
    \rule{5cm}{0.4pt} \\
    Md. Imtiaj Alam Sajin \\
\end{tabular}

\newpage

%---------------------------------------------------------------------
% ABSTRACT
%---------------------------------------------------------------------
\section*{Abstract}
\addcontentsline{toc}{section}{Abstract}
Large language models now translate natural language questions into SQL
with high execution accuracy, and are increasingly proposed as database
interfaces for non-expert users. That accuracy is measured almost
entirely as whether the returned rows are correct. Whether the generated
query is an \emph{efficient} way to obtain them is largely unexamined: a
semantically correct query can be arbitrarily slower than an equivalent
formulation, and the model knows nothing of the physical design, the
data distribution, or the cost of the plan its text will produce. This
proposal argues that the knowledge needed to close that gap already
exists inside every relational system, in the optimiser's cost model,
and can be consulted before a query is executed. I propose a
generate-and-rank architecture in which the model produces several
semantically-equivalent formulations of one question, each is costed
through \texttt{EXPLAIN} without execution, and the cheapest is
selected. The obvious objection is that this assumes the cost model
ranks alternatives correctly. My own completed measurement study,
summarised here as preliminary evidence, shows that assumption fails in
identifiable regimes: across 23{,}480 measurements on two independently
written engines, 98.2\% of the cases in which PostgreSQL chose a slower
access path than one available to it carried an essentially correct row
estimate, placing the fault in the cost model rather than in cardinality
estimation. The proposal therefore includes a reliability guard that
suppresses cost-based selection where the model is known to misrank, and
an evaluation plan built on execution time rather than accuracy alone,
using Spider, BIRD and BIRD's Valid Efficiency Score.

\vspace{0.5em}
\noindent
\textbf{Keywords:} text-to-SQL, large language models, query
optimisation, cost models, cardinality estimation, query efficiency.

\newpage
\tableofcontents
\newpage

%=====================================================================
\section{Introduction}
\label{sec:intro}
%=====================================================================

\subsection{Motivation}

The ability to ask a database a question in ordinary language has been a
research goal for decades, and it has become a practical one only
recently. Sequence-to-sequence models first made the task tractable at
scale~\cite{zhong2017seq2sql}, schema-aware encoders improved it
substantially~\cite{wang2020ratsql}, and large language models have now
pushed execution accuracy on the standard cross-domain benchmark past
levels that seemed distant a few years
ago~\cite{yu2018spider, pourreza2023dinsql, gao2024dailsql}. Systems
built on this capability are being offered as database interfaces for
users who do not write SQL.

That framing carries an assumption worth examining. When a database
administrator writes a query, correctness is the beginning of the job,
not the end of it. Two queries can return byte-identical results and
differ by orders of magnitude in the time and I/O they consume, because
SQL is a declarative language and the same request can be expressed in
forms that lead the optimiser to very different plans. A subquery that
could have been a join, a predicate written so that it cannot use an
index, an unnecessary \texttt{DISTINCT} that forces a sort: none of
these change the answer, and all of them change the cost.

A language model generating SQL has no access to the information needed
to avoid these. It sees a question and a schema. It does not see which
columns are indexed, how many rows the table holds, how values are
distributed, whether the data is physically clustered, or how much
memory the buffer pool has. Those are exactly the inputs a query
optimiser consumes. The model is being asked to make a decision with
consequences it cannot observe.

\subsection{Problem statement}
\label{sec:problem}

The evaluation methodology of the field reflects this gap. The dominant
metric is execution accuracy, the fraction of generated queries whose
result set matches the reference. Under that metric a query that takes
four minutes and a query that takes forty milliseconds score
identically, provided both are correct. The BIRD
benchmark~\cite{li2023bird} is the notable exception and deserves
credit for it: alongside execution accuracy it introduced the Valid
Efficiency Score, which weights correct queries by their runtime
relative to the reference. Its authors observed that generated queries
are frequently far slower than the human-written ones. That observation
has not yet reshaped how systems are built. Efficiency remains a
reported number rather than an objective the generation process
optimises for.

The problem this proposal addresses is therefore narrow and concrete:

\begin{quote}
\emph{Text-to-SQL systems generate queries without any knowledge of what
those queries will cost to execute, and are evaluated under metrics that
do not penalise the difference. The information required to close that
gap is already computed by the database, in the optimiser's cost model,
and is available before execution. It is not currently used.}
\end{quote}

\subsection{Proposed direction}

The direction I propose is to place the optimiser inside the generation
loop rather than after it. Instead of producing one query and executing
it, the system produces several semantically-equivalent formulations of
the same question, submits each to the database's \texttt{EXPLAIN}
facility, which returns an estimated cost without executing anything,
and selects the formulation the optimiser expects to be cheapest.

This is a modest architectural idea, and its value depends entirely on a
question that the text-to-SQL literature has had no reason to ask:
\emph{is the optimiser's cost estimate a trustworthy ranking signal?}
The database systems literature has asked it, and the answer is
qualified~\cite{leis2015good, wu2013predicting}. My own completed
measurement study, described in Section~\ref{sec:prelim}, answers it
more sharply for the specific decision that matters here and finds that
the estimate is reliable in most regimes and systematically misleading
in a few that can be characterised in advance.

That finding does not defeat the proposal. It specifies it. A
cost-aware system that consults the optimiser naively will inherit the
optimiser's blind spots; one that knows where those blind spots lie can
decline to act on the signal there. Building that distinction into the
design is the substance of what I am proposing.

\subsection{Research questions}
\label{sec:rqs}

\begin{description}[leftmargin=2.1cm, style=nextline, font=\bfseries]
\item[RQ1] How large is the efficiency gap between LLM-generated SQL and
expert-written reference SQL on established benchmarks, when both are
measured by execution time on identical physical designs rather than by
result correctness?

\item[RQ2] Do semantically-equivalent reformulations of a single natural
language question exhibit enough variation in execution cost for
selection among them to be worthwhile, and how large is that spread?

\item[RQ3] Can the optimiser's estimated cost, obtained without
execution, rank those candidate formulations well enough to serve as a
selection signal?

\item[RQ4] In which regimes does that ranking fail, and can those regimes
be detected cheaply enough at generation time to suppress cost-based
selection where the signal is unreliable?

\item[RQ5] What does the resulting system cost in latency, and is the
efficiency gained worth the additional optimiser calls and generated
candidates?
\end{description}

\subsection{Contributions}
\label{sec:contrib}

Should the proposed work be carried out, its intended contributions are
an efficiency-first characterisation of the text-to-SQL efficiency gap
under controlled physical design; a candidate-generation and cost-ranking
architecture that uses the optimiser as a pre-execution oracle; a
reliability guard derived from measured optimiser failure regimes rather
than assumed; and a public artifact permitting reproduction. The
proposal is deliberately scoped so that RQ1 and RQ2 are answerable
independently of whether the full system succeeds, and a negative
result on RQ3 would itself be a reportable finding about the limits of
cost-model reuse.

%=====================================================================
\section{Background}
\label{sec:background}
%=====================================================================

\subsection{The text-to-SQL pipeline}

Contemporary systems decompose the task into roughly three stages.
\emph{Schema linking} identifies which tables and columns the question
refers to, a step that dominates error rates on large schemas and has
received corresponding attention~\cite{li2023resdsql, talaei2024chess}.
\emph{Generation} produces candidate SQL, either by constrained decoding
that enforces syntactic and schema validity during
generation~\cite{scholak2021picard}, or by prompting a large model with
in-context examples~\cite{pourreza2023dinsql, gao2024dailsql}.
\emph{Verification} checks the result, typically by executing the query
and repairing it if it errors, and increasingly by generating several
candidates and choosing among them~\cite{pourreza2024chase}.

That last development matters for this proposal. Multi-candidate
generation is already established practice; systems produce a set of
plausible queries and select one. What they select \emph{on} is
correctness: a model-based preference score, self-consistency across
samples, or agreement of executed results. Nothing in the selection step
considers cost. The infrastructure this proposal needs, therefore,
largely exists; the criterion it applies is different.

\subsection{Cost-based query optimisation}

When a relational engine receives SQL, it enumerates alternative
execution plans and estimates the cost of each, selecting the cheapest.
The framework dates to System~R~\cite{selinger1979access} and remains
essentially intact in modern engines~\cite{chaudhuri1998overview,
graefe1993query}. Cost is computed from cardinality estimates, how many
rows each operator will produce, derived from stored statistics, and
from a set of tunable constants expressing the relative expense of
sequential and random I/O and of CPU work.

Two properties of this machinery matter here. First, it is available
without execution: \texttt{EXPLAIN} returns the chosen plan and its
estimated cost in microseconds, having touched no data. That is what
makes pre-execution ranking feasible at all. Second, its accuracy is
imperfect and unevenly distributed. Cardinality errors compound through
join trees~\cite{ioannidis1991propagation}, and the seminal empirical
assessment of optimiser quality found misestimation to be the dominant
source of poor join plans~\cite{leis2015good, leis2018query}. Cost
models have separately been found usable for runtime prediction only
after calibration~\cite{wu2013predicting}.

\subsection{Why the two literatures have not met}

Text-to-SQL research treats the database as a black box that returns
rows, and evaluates the text that goes into it. Query optimisation
research treats the SQL as given, and works on what happens after it
arrives. The interface between them, the fact that the choice of SQL
text materially constrains the plans the optimiser can consider, falls
between the two.

There are early signs of contact. CodexDB~\cite{trummer2022codexdb}
generates query-processing code from natural language instructions.
Recent work uses language models to apply rewrite
rules~\cite{zhou2024llmr2} and asks directly whether they can serve as
query optimisers~\cite{zhu2024llmquery}. These operate on SQL that
already exists. The proposal here concerns the moment before that, when
the SQL is still being chosen.

%=====================================================================
\section{Related Work}
\label{sec:related}
%=====================================================================

\paragraph{Text-to-SQL systems and benchmarks.}
Spider~\cite{yu2018spider} established cross-domain evaluation, where
test databases are unseen at training time, and remains the standard
reference. BIRD~\cite{li2023bird} moved to larger, dirtier, more
realistic databases and, importantly for this proposal, introduced the
Valid Efficiency Score to capture that a correct query may still be a
bad one. Surveys document the transition from sequence-to-sequence
models to large language
models~\cite{katsogiannis2023survey}. Representative recent systems span
decomposed prompting~\cite{pourreza2023dinsql}, systematic prompt
engineering~\cite{gao2024dailsql}, open-source specialised
models~\cite{li2024codes}, multi-agent
decomposition~\cite{wang2025macsql} and candidate selection under
multi-path reasoning~\cite{pourreza2024chase}. The through-line is that
selection among candidates is now standard and is performed on
correctness signals.

\paragraph{Efficiency of generated queries.}
This is the thin part of the literature, which is why it is the opening.
BIRD reports that generated queries are often markedly slower than
reference queries, and its Valid Efficiency Score quantifies this, but
the metric is generally reported as a secondary column rather than
optimised against. I am not aware of a system that treats expected
execution cost as a first-class objective during generation, and
establishing whether one exists is the first task of the proposed work.

\paragraph{Learned and hybrid optimisation.}
A parallel line of work replaces or augments the analytic cost model
with a learned one~\cite{marcus2019neo, marcus2021bao}, motivated by
the premise that the shipped model is miscalibrated enough to be worth
replacing. Related work trains cardinality estimators against a
plan-cost objective rather than an accuracy objective, precisely because
the two diverge~\cite{negi2021flow}. The proposal here is complementary
and deliberately more conservative: rather than replacing the cost
model, it reuses the existing one as a comparator among queries the
system itself authored, and treats the model's known weaknesses as
constraints on when to trust it.

\paragraph{Cost model reliability.}
Whether optimiser cost predicts runtime has been studied
directly~\cite{wu2013predicting} and comprehensively for cardinality
estimation~\cite{han2021cardinality}. Surveys note the scarcity of
controlled studies isolating individual optimiser
components~\cite{lan2021survey}. Section~\ref{sec:prelim} contributes a
measurement of exactly this kind for the single-relation access-path
decision, which is the decision most affected by how a predicate is
written and therefore the one most relevant to query formulation.

%=====================================================================
\section{Preliminary Evidence}
\label{sec:prelim}
%=====================================================================

This section reports measured results. They come from a completed
experimental study I conducted on PostgreSQL~17.1 and MariaDB~10.4.28,
comprising 23{,}480 measurements across seven table scales, eleven
controlled dataset configurations and ten index configurations, with
every dataset generated from a fixed seed so that the true row count of
every predicate is known rather than estimated. They are included here
because they determine whether the architecture proposed in
Section~\ref{sec:approach} is sound, and because they were the reason I
arrived at that architecture.

\subsection{The cost model misranks even when estimates are correct}

The accepted account of poor plan selection attributes it to
cardinality misestimation. Testing that on the access-path decision
produced the opposite result. Among the cases where the optimiser chose
a demonstrably slower access path than one available to it, the median
estimation error was 1.017, where a perfect estimate is 1.0, and 98.2\%
of such cases carried an essentially correct row estimate
(Table~\ref{tab:dissoc}).

\begin{table}[htbp]
\centering
\caption{Estimation quality among PostgreSQL's misselected access
paths. The optimiser knew the row counts and ranked the alternatives
wrongly regardless.}
\label{tab:dissoc}
\small
\begin{tabular}{@{}lr@{}}
\toprule
\textbf{Measure} & \textbf{Value} \\
\midrule
Misselected cases examined & 113 \\
Median $q$-error among them & 1.017 \\
Share with an essentially correct estimate & 98.2\% \\
A perfect estimate, for reference & 1.000 \\
\bottomrule
\end{tabular}
\end{table}

The implication for this proposal cuts both ways. On the one
hand, it means that improving the statistics, which is what most tuning
advice targets, cannot repair these rankings, so a system that wants
better decisions must intervene elsewhere, which is what selecting among
formulations does. On the other hand, it means the cost model itself is
the faulty component, and a design that trusts it uncritically is
building on the very thing measured to be unreliable.

\subsection{A characterised failure regime}

The largest effect found was specific and mechanically explicable. When
a BRIN index is present alongside a B-tree on the same column, a
configuration widely recommended on the grounds that BRIN indexes are
small and cheap, PostgreSQL's misselection rate on the affected
workload rose from 2.3\% to 73.3\%, with a worst case of 194$\times$ on
a controlled test in which the only change was the presence of the BRIN
index.

The cause was traced to the source. In \texttt{choose\_bitmap\_and()},
when two indexes can answer the same predicate, the optimiser retains
whichever is cheaper to scan, judged by the cost of reading the index
alone and excluding the heap fetches the query must then perform. A BRIN
index is small by design and therefore wins that comparison almost
always (Table~\ref{tab:brincause}).

\begin{table}[htbp]
\centering
\caption{The deciding comparison in a measured case. The index retained
is the one that runs 194$\times$ slower.}
\label{tab:brincause}
\small
\begin{tabular}{@{}lrr@{}}
\toprule
\textbf{Index} & \textbf{Cost used to decide} & \textbf{Actual query time} \\
\midrule
BRIN & 12.13 \quad (retained) & 65.95 ms \\
B-tree & 213.35 \quad (discarded) & 0.34 ms \\
\bottomrule
\end{tabular}
\end{table}

This matters to the proposal for a reason beyond the specific defect:
the regime is \emph{detectable}. It requires two indexes of differing
kinds on one column, and it was confined to tables still resident in the
buffer pool. Both conditions are readable from the catalogue at
generation time. A failure regime that can be recognised in advance can
be guarded against, and that is the mechanism
Section~\ref{sec:guard} proposes.

\subsection{A shared weakness across two engines}

Correlated conjunctive predicates were mishandled by both systems.
PostgreSQL's estimation error rose with correlation to 98.65, while
MariaDB estimated the same predicates perfectly, at exactly 1.00, until
no single index covered both columns, at which point its error jumped to
55.84. Two independently written engines failing at the same structural
point indicates a shared design assumption rather than a defect in
either product: both compose per-column selectivities by multiplication
when no single index can measure the combination directly.

Predicate structure is precisely what a text-to-SQL system controls. If
composition across columns is where estimation reliably breaks, then how
a question's filters are distributed across predicates is a lever the
generator holds, and this is a second reason to expect the choice of
formulation to matter.

\subsection{What the preliminary evidence establishes}

Three things, which together define the proposal's shape. First, the
premise is real: cost-based ranking is fallible in ways not attributable
to statistics, so the proposal must be built with that in mind rather
than assuming an oracle. Second, the fallibility is bounded and
characterisable, not diffuse, so a guard is feasible. Third, predicate
formulation interacts with the failure modes, so selecting among
formulations is a plausible lever rather than a speculative one.

%=====================================================================
\section{Proposed Approach}
\label{sec:approach}
%=====================================================================

\subsection{Architecture}

The proposed system inserts two stages between generation and execution.

\begin{enumerate}[leftmargin=*, itemsep=0.35em]
\item \textbf{Schema and physical-design context.} Beyond table and
column names, the prompt is supplied with the physical design: which
columns carry indexes and of what kind, approximate table cardinalities,
and the presence of extended statistics. This is read from the
catalogue and costs one query per database.

\item \textbf{Multi-candidate generation.} The model produces $k$
candidate SQL statements for one question, prompted for
\emph{formulational} diversity: join order and join form, subquery
versus \texttt{JOIN} versus \texttt{EXISTS}, predicate placement,
explicit projection instead of \texttt{SELECT *}.

\item \textbf{Semantic equivalence filtering.} Candidates that do not
agree on results are not comparable, so this step is a precondition
rather than an optimisation. On the benchmark databases, agreement is
checked by execution on the provided instance; the limits of this are
discussed in Section~\ref{sec:risks}.

\item \textbf{Cost-based ranking.} Each surviving candidate is submitted
under \texttt{EXPLAIN}, without \texttt{ANALYZE}, so no execution
occurs. The optimiser returns its chosen plan and estimated total cost.

\item \textbf{Reliability guard.} Before the ranking is acted upon, the
plans are inspected for the regimes in which cost is known to misrank.
Where a regime is detected, cost-based selection is suppressed and the
system falls back to the model's own first-ranked candidate.

\item \textbf{Execution.} The selected query is executed and returned.
\end{enumerate}

\subsection{The reliability guard}
\label{sec:guard}

This is the component that distinguishes the proposal from the naive
design, and it is derived from measurement rather than intuition. Three
conditions are proposed for the initial guard, each drawn from
Section~\ref{sec:prelim}:

\begin{description}[leftmargin=1.6cm, style=nextline, font=\itshape]
\item[Co-located heterogeneous indexes] If the plan's relation carries
two indexes of different kinds on the same column, the measured
misranking regime applies and the cost comparison between index paths is
not to be trusted.

\item[Correlated conjunctive predicates] If the plan applies a
conjunction over two or more columns of the same relation with no single
index covering the combination, estimates are composed multiplicatively
and were measured to be unreliable in both engines.

\item[Narrow cost margins] If the estimated costs of the top candidates
lie within a threshold of one another, the ranking is treated as a tie
and the model's preference is retained. The threshold is to be
calibrated empirically rather than assumed, and calibrating it is an
experiment in its own right.
\end{description}

The guard is intentionally conservative. Its purpose is not to improve
rankings but to avoid acting on bad ones, so its failure mode is
reverting to current practice rather than performing worse than it.

\subsection{Why not simply execute the candidates?}

Executing every candidate and keeping the fastest would give a better
ranking than any estimate. It is also self-defeating: the cost of the
selection exceeds the cost of the query, and on an expensive query the
system would have to run the slow formulations to discover they are
slow. Pre-execution estimation is what makes the idea deployable, and it
is why cost-model reliability is the crux rather than a detail.

%=====================================================================
\section{Evaluation Plan}
\label{sec:eval}
%=====================================================================

\subsection{Datasets and physical design}

Spider~\cite{yu2018spider} for cross-domain breadth and comparability,
and BIRD~\cite{li2023bird} as the primary target because its databases
are large enough for physical design to matter and because it already
supplies an efficiency metric. A limitation of both must be stated
plainly: their databases ship with minimal indexing, so an efficiency
study on them as distributed would measure a setting no production
system resembles. The proposed work therefore evaluates each database
twice, once as shipped and once under a defined physical design applied
uniformly, with the index set published as part of the artifact.

\subsection{Metrics}

\begin{description}[leftmargin=2.4cm, style=nextline, font=\bfseries]
\item[Execution accuracy] Result-set equivalence against the reference.
This is a gate, not a target: any efficiency gain accompanied by an
accuracy loss is a failure of the method.

\item[Valid Efficiency Score] BIRD's efficiency measure, for
comparability with published work.

\item[Median and tail latency ratio] Execution time of the selected
query over that of the reference query, reported at the median and at
the 95th percentile. The tail is reported separately because the
motivating pathology is rare and severe rather than typical.

\item[Selection regret] Execution time of the selected candidate over
that of the fastest candidate generated. This isolates the quality of
the ranking decision from the quality of the candidate set, and is the
metric that answers RQ3 directly.

\item[Overhead] Added wall-clock time from extra generation and
\texttt{EXPLAIN} calls, reported per query.
\end{description}

\subsection{Baselines and ablations}

Baselines: the reference SQL as an upper bound; single-candidate
generation as current practice; and random selection among candidates,
which separates a genuine ranking effect from the effect of merely
generating more candidates. Ablations: cost ranking without the guard,
to measure what the guard contributes; the guard's three conditions
individually; schema-only prompting against physical-design-aware
prompting, to determine whether the context alone changes what the model
writes; and $k$ swept to establish where candidate count stops paying.

\subsection{Answering the research questions}

RQ1 is answered by the baseline comparison of reference against
single-candidate generation. RQ2 by the spread of execution times within
each candidate set, which requires executing all candidates during
evaluation though not in deployment. RQ3 by selection regret. RQ4 by
partitioning selection regret across the guard conditions. RQ5 by
overhead measured against latency improvement.

\subsection{Measurement protocol}

The protocol from my completed study transfers directly and is adopted:
repeated execution with the first run discarded as cache-warming, the
median of the remainder reported, dispersion reported alongside it,
parallelism and just-in-time compilation disabled so they cannot
confound the comparison, and a fixed buffer pool. Every raw measurement
is retained rather than only the summary, so that dispersion claims can
be checked independently.

%=====================================================================
\section{Risks and Limitations}
\label{sec:risks}
%=====================================================================

\paragraph{Semantic equivalence cannot be fully verified.}
Agreement of results on one database instance is evidence of
equivalence, not proof; two queries can agree on the data present and
diverge on data absent. This is a real and unavoidable limitation. The
mitigation is to constrain candidate generation to a set of
transformations that are equivalence-preserving by construction, to
report agreement checking as a filter rather than a guarantee, and to
treat any accuracy regression as disqualifying.

\paragraph{The cost model may not rank formulations well enough.}
RQ3 may return a negative answer. This is the principal scientific risk
and it is accepted deliberately: the preliminary evidence shows the
model misranking \emph{access paths} within one query, which is related
to but not identical with ranking across query formulations. A negative
result would be reportable and would bound a design others might
otherwise attempt.

\paragraph{Benchmark databases may be too small.}
If the databases are small enough that every plan is fast, formulation
will not matter and no method can demonstrate an effect. The scale
boundary measured in my completed study, where behaviour changed sharply
around buffer-pool residency, indicates that this is a real hazard and
suggests scaling the databases as a contingency.

\paragraph{Results will be engine-specific.}
Cost models differ. The completed study found selection quality to
differ sharply between PostgreSQL and MariaDB, so any finding here
should be expected to be engine-dependent until shown otherwise, and
the evaluation is planned on PostgreSQL with a smaller replication on a
second engine for exactly that reason.

\paragraph{Cost and access.} Multi-candidate generation multiplies
inference cost by $k$. Open-weight models~\cite{li2024codes} are the
mitigation if API cost becomes limiting.

%=====================================================================
\section{Work Plan}
\label{sec:plan}
%=====================================================================

\begin{table}[htbp]
\centering
\caption{Indicative phasing. Phases~1 and~2 are self-contained and
produce reportable findings regardless of the outcome of later phases.}
\label{tab:timeline}
\small
\begin{tabular}{@{}clp{7.4cm}@{}}
\toprule
\textbf{Phase} & \textbf{Duration} & \textbf{Work} \\
\midrule
1 & Weeks 1--3 & Harness, benchmark loading, physical designs,
measurement protocol. Establish the efficiency gap (RQ1). \\
2 & Weeks 4--6 & Multi-candidate generation and equivalence filtering.
Measure the cost spread within candidate sets (RQ2). \\
3 & Weeks 7--9 & Cost-based ranking; measure selection regret (RQ3). \\
4 & Weeks 10--12 & Guard conditions, calibration, ablations (RQ4). \\
5 & Weeks 13--15 & Overhead, second-engine replication (RQ5). \\
6 & Weeks 16--18 & Write-up and artifact preparation. \\
\bottomrule
\end{tabular}
\end{table}

%=====================================================================
\section{Expected Outcomes and Conclusion}
\label{sec:conclusion}
%=====================================================================

The anticipated outcome, stated as an expectation rather than a result,
is that a measurable efficiency gap exists between generated and
reference SQL; that semantically-equivalent formulations of one question
vary enough in cost for selection to be worthwhile; and that optimiser
cost ranks them usefully in most regimes while failing in the specific
ones the preliminary study identified, so that a guarded design
outperforms both an unguarded one and current single-candidate practice.

The value of the proposal does not depend on all of these holding. RQ1
and RQ2 are descriptive and will yield findings either way. A negative
answer on RQ3 would be a useful result about the limits of reusing an
internal cost model as an external ranking signal, and it would be
grounded in measurement rather than assumption.

What the proposal contributes conceptually is a change of objective.
Text-to-SQL is currently posed as a translation problem, scored on
whether the translation is faithful. A database administrator asked to
write the same query treats faithfulness as the entry requirement and
spends their effort on what the query will cost. Building systems that
are evaluated the way a practitioner would evaluate them, and giving
them access to the cost information the database already computes, seems
a reasonable next step, and the preliminary evidence gathered so far
indicates both that the information is worth using and that it must be
used carefully.

\vspace{0.8em}
\noindent
\textbf{Relationship to my completed coursework.} The measurement study
summarised in Section~\ref{sec:prelim} was undertaken in this course and
is reported in full separately. It began as groundwork for this
proposal, to establish whether the optimiser's cost model was
trustworthy enough to build upon, and grew into a self-contained
empirical study when the answer turned out to be more interesting than
expected. The two documents are complementary: one measures the
component, the other proposes what to build with it.

%---------------------------------------------------------------------
\newpage
\printbibliography[heading=bibnumbered, title={References}]

\end{document}
